\chapter{Core Lemmas}

\section{Chain Rule Lemma}

The chain rule is the elementary engine behind the capacity calculus.

\begin{lemma}[Transcript chain rule]
For transcript variables \(T_1,\ldots,T_T\),
\[
I(X;T_1,\ldots,T_T)
=
\sum_{t=1}^{T} I(X;T_t \mid T_1,\ldots,T_{t-1}) .
\]
\end{lemma}

\begin{proof}
This is the standard chain rule for mutual information.
\end{proof}

\section{Non-Amplification Lemma}

\begin{lemma}[Entropy non-amplification]
If each refinement step has conditional information gain at most \(C_t\), then the transcript cannot expose more than \(\sum_t C_t\) total information about \(X\).
\end{lemma}

\begin{proof}
Apply the transcript chain rule and the per-step bounds.
\end{proof}

\section{Terminal Obstruction Lemma}

\begin{lemma}[Terminal obstruction]
If a target requires entropy demand \(D\) and a normalized refinement sequence has total capacity less than \(D\), then the target cannot be fully resolved by that sequence.
\end{lemma}

\begin{proof}
Full resolution would require transcript information at least \(D\), contradicting non-amplification.
\end{proof}
